\documentclass[
  fontsize=12,
  DIV=15,
]{scrreprt}

%% Задаем размер полей путем указания области, занимаемой текстом
% \usepackage[a4paper, total={7in, 9in}]{geometry}

%~~~~~~~~~~~~~~~~~~~~~~~~~~~~~~~~~~~~~~~~~~~~~~~~~~~~~~~~~~~~~~~~~~
%              Настройка шрифтов
%~~~~~~~~~~~~~~~~~~~~~~~~~~~~~~~~~~~~~~~~~~~~~~~~~~~~~~~~~~~~~~~~~~
\usepackage[no-math]{fontspec}
% При использовании LuaLaTeX отключаем следующие два пакета
% \usepackage{xunicode}
% \usepackage{xltxtra}



\usepackage{polyglossia}
\setdefaultlanguage[indentfirst=true,spelling=modern]{russian}
\setotherlanguage{english}

\usepackage{amsmath}
\usepackage{unicode-math}

\unimathsetup{
  mathbf=sym,
  mathrm=sym,
  % math-style=TeX,
  % bold-style=TeX
}

% \setmainfont{CMU Serif}
% \setromanfont{CMU Serif}
% \setsansfont{CMU Sans Serif}
% \setmonofont{CMU Typewriter Text}


\setmainfont{IBM Plex Serif}
\setromanfont{IBM Plex Serif}
\setsansfont{IBM Plex Sans}
\setmonofont{IBM Plex Mono}
\setmonofont{JuliaMono}

% Шрифты математических формул
\setmathfont{Latin Modern Math}
% \setmathfont{TeX Gyre Bonum Math}
% \setmathfont{TeX Gyre Pagella Math}
% \setmathfont{TeX Gyre Schola Math}
% \setmathfont{TeX Gyre Termes Math}
% \setmainfont{Libertinus Math}
% \setmainfont{Cambria Math}
% \setmathfont{XITS Math}
% \setmathfont{STIX}
% \setmathfont{DejaVu Math TeX Gyre} % нет
% \setmathfont{Asana-Math fonts} % нет
% \setmainfont{Lucida Bright Math} % нет
% \setmainfont{Fira Math} % нет
% \setmainfont{Minion Math} % нет
%~~~~~~~~~~~~~~~~~~~~~~~~~~~~~~~~~~~~~~~~~~~~~~~~~~~~~~~~~~~~~~~~~~~~~~~~~~~~~~~
%								Пакеты для математики
%~~~~~~~~~~~~~~~~~~~~~~~~~~~~~~~~~~~~~~~~~~~~~~~~~~~~~~~~~~~~~~~~~~~~~~~~~~~~~~~
% Некоторые дополнительные символы (например стрелки)
% \usepackage{latexsym}
% Рукописное начертание (с завитушками)
% \usepackage{mathrsfs}
% Теоремы и т.п.
% \usepackage{amsthm}
% Пакет для объявления новых математических операторов
% \usepackage{amsopn}
% \usepackage{tensor}
\usepackage{physics}
%~~~~~~~~~~~~~~~~~~~~~~~~~~~~~~~~~~~~~~~~~~~~~~~~~~~~~~~~~~~~~~~~~~
%       Настройка пакета hyperref
%~~~~~~~~~~~~~~~~~~~~~~~~~~~~~~~~~~~~~~~~~~~~~~~~~~~~~~~~~~~~~~~~~~
\usepackage{hyperref}
\hypersetup{
  colorlinks,
  citecolor=black,
  filecolor=black,
  linkcolor=black,
  urlcolor=black,
  bookmarksdepth=section, 
  unicode=true,       
  pdftoolbar=true,       
  pdfmenubar=true,    
  pdffitwindow=false,  
  pdfstartview={FitH},  
  pdftitle={Название},
  pdfauthor={Автор}
}
%~~~~~~~~~~~~~~~~~~~~~~~~~~~~~~~~~~~~~~~~~~~~~~~~~~~~~~~~~~~~~~~~~~
%                Пакет для листингов
%~~~~~~~~~~~~~~~~~~~~~~~~~~~~~~~~~~~~~~~~~~~~~~~~~~~~~~~~~~~~~~~~~~
% Пакет для настройки verbatim'а
% \usepackage{fancyvrb}
% Пакет для листингов. Опция chapter указывает на нумерацию листингов по главам
\usepackage{minted}
% friendly monokai manni borland trac paraiso-light fruity emacs autumn
% \usemintedstyle{bw} % черно-белый
% \usemintedstyle{friendly}
\setminted{
  autogobble=true,
  breaklines=true,
  breakanywhere=true,
  obeytabs=false,
  tabsize=2,
  linenos=true,
  %%stepnumber=5
}
% Пакеты для псевдокода
% \usepackage{verbatim}% Для исходного кода, цитат и т.д. и т.п.
% Для алгоритмов
% \usepackage{algorithm}
% \usepackage{algpseudocode}
% \usepackage{algorithmicx}
%~~~~~~~~~~~~~~~~~~~~~~~~~~~~~~~~~~~~~~~~~~~~~~~~~~~~~~~~~~~~~~~~~~
%              Графика в TeX
%~~~~~~~~~~~~~~~~~~~~~~~~~~~~~~~~~~~~~~~~~~~~~~~~~~~~~~~~~~~~~~~~~~
\usepackage{graphicx} % графика
\usepackage{xcolor} % цвет
%~~~~~~~~~~~~~~~~~~~~~~~~~~~~~~~~~~~~~~~~~~~~~~~~~~~~~~~~~~~~~~~~~~~~~~~~~~~~~~~
%          BibLaTeX
%~~~~~~~~~~~~~~~~~~~~~~~~~~~~~~~~~~~~~~~~~~~~~~~~~~~~~~~~~~~~~~~~~~~~~~~~~~~~~~~
\usepackage[autostyle]{csquotes}
\usepackage{url}

\usepackage[%
  backend=biber,
  bibstyle=gost-numeric,
  defernumbers=true,
  % не менять местами заголовок и список авторов, если авторов больше четырех
  movenames=false,
  % не использовать сокращение списка авторов
  maxbibnames=99,
  sorting=ydnt
]{biblatex}
% Файл с литературой
\bibliography{lit}

%~~~~~~~~~~~~~~~~~~~~~~~~~~~~~~~~~~~~~~
%   Определение новых команд
%~~~~~~~~~~~~~~~~~~~~~~~~~~~~~~~~~~~~~~
% Стандартные множества
% \newcommand{\setN}{\mathbb{N}}
% \newcommand{\setZ}{\mathbb{Z}}
% \newcommand{\setQ}{\mathbb{Q}}
% \newcommand{\setR}{\mathbb{R}}
% \newcommand{\setC}{\mathbb{C}}

% \renewcommand{\det}{\mathrm{det}}% Определитель
% \renewcommand{\d}{\mathrm{d}}% Инфинитозимальный дифференциал
